% ========== 2.2 SYSTEM STAKEHOLDERS ==========
% Identification of users and their roles

\section{2.1 System Stakeholders}
\addcontentsline{toc}{section}{2.1 System Stakeholders}
\label{sec:system-stakeholders}

\lettrine{T}{he Symphony system} serves multiple stakeholder groups with distinct needs and expectations. Understanding these stakeholders is essential for defining appropriate system requirements and ensuring the solution addresses real-world needs across different usage contexts.

Individual developers represent the primary user group, including professional programmers working on personal projects, freelance developers serving multiple clients, and students learning software development. These users require tools that enhance individual productivity, provide intelligent assistance without overwhelming complexity, and support rapid learning of new technologies and frameworks. Their success depends on the system's ability to amplify individual capabilities while maintaining simplicity and ease of use.

Development teams form another critical stakeholder group, ranging from small teams of three to ten developers to large organizations with hundreds of programmers. Team stakeholders need capabilities that facilitate collaboration, maintain code consistency across team members, and enable knowledge sharing. Technical leaders and project managers within these teams require visibility into development activities, tools for enforcing standards, and mechanisms for tracking progress and quality metrics.

Educational institutions and training organizations represent stakeholders focused on teaching and learning. Computer science educators need systems that support curriculum delivery, enable student progress tracking, and provide appropriate scaffolding for learners at different skill levels. Students benefit from tools that explain concepts clearly, tolerate mistakes during learning, and provide constructive feedback that accelerates skill development.

Enterprise organizations evaluating the system for adoption have comprehensive requirements spanning technical capabilities, security compliance, and business considerations. Information technology departments need deployment flexibility and integration with existing infrastructure. Security teams require robust access controls, audit capabilities, and compliance with regulatory standards. These stakeholders prioritize reliability, maintainability, and long-term vendor support over cutting-edge features.
